\documentclass[11pt]{article}
\usepackage[margin=1in]{geometry}
\usepackage[T1]{fontenc}
\usepackage{lmodern}
\usepackage{amsmath}
\usepackage{booktabs}
\usepackage{array}
\usepackage{cleveref}

% Significance stars rendered as a superscript, the estout/esttab convention.
\newcommand{\sym}[1]{\textsuperscript{#1}}

\title{A Spanned Booktabs Regression Table}
\author{Test Corpus}
\date{}

\begin{document}
\maketitle

\section{Introduction}

This fixture exercises the structural features that distinguish a
publication-grade regression table from a plain grid. The table in
\cref{tab:spanned} reports four specifications grouped into two blocks. A first
header level uses \verb|\multicolumn| to place two group titles, ``Baseline''
and ``Robustness,'' each spanning a pair of estimate columns; a second header
level names the individual specifications beneath them. Two partial rules drawn
with \verb|\cmidrule(lr)| underline the spanners without running the full width
of the table, which is the booktabs idiom for showing that the group titles
govern only their own columns.

Below the header, each coefficient is followed on the next physical row by its
standard error in parentheses, and \verb|\addlinespace| separates the
coefficient blocks so the eye groups each estimate with its error. Significance
is marked with superscript stars produced by a \verb|\sym| macro, and a note
line beneath the rule records the estimator and the star thresholds. The point
of the fixture is that all of this carries meaning: the column grouping, the
partial rules, the coefficient/standard-error pairing, and the note must survive
conversion for the table to remain readable.

\section{Results}

\Cref{tab:spanned} summarizes the four models. The interaction term
$\text{Treatment}\times\text{Post}$ is the estimand of interest, and its
magnitude is stable across the baseline and robustness blocks, with
$\widehat{\beta}\approx 0.11$ throughout.

\begin{table}[htbp]
\centering
\caption{Difference-in-differences estimates across baseline and robustness
  specifications.}
\label{tab:spanned}
\begin{tabular}{l cc cc}
\toprule
& \multicolumn{2}{c}{Baseline} & \multicolumn{2}{c}{Robustness} \\
\cmidrule(lr){2-3}\cmidrule(lr){4-5}
& (1) & (2) & (3) & (4) \\
& OLS & FE & FE+Cov. & IV \\
\midrule
Treatment            & 0.184\sym{***} & 0.171\sym{***} & 0.166\sym{***} & 0.158\sym{**}  \\
                     & (0.041)        & (0.039)        & (0.038)        & (0.062)        \\
\addlinespace
Post period          & 0.072\sym{**}  & 0.064\sym{**}  & 0.059\sym{*}   & 0.061\sym{*}   \\
                     & (0.031)        & (0.030)        & (0.030)        & (0.033)        \\
\addlinespace
Treatment $\times$ Post & 0.119\sym{***} & 0.111\sym{***} & 0.108\sym{***} & 0.114\sym{***} \\
                     & (0.028)        & (0.027)        & (0.027)        & (0.035)        \\
\addlinespace
Covariate index      &                &                & -0.031         & -0.028         \\
                     &                &                & (0.024)        & (0.025)        \\
\addlinespace
Constant             & 1.402\sym{***} & 1.388\sym{***} & 1.451\sym{***} & 1.460\sym{***} \\
                     & (0.066)        & (0.061)        & (0.070)        & (0.072)        \\
\midrule
Unit fixed effects   & No  & Yes & Yes & Yes \\
Observations         & 4{,}812 & 4{,}812 & 4{,}812 & 4{,}517 \\
$R^{2}$              & 0.214 & 0.286 & 0.301 & 0.255 \\
\bottomrule
\addlinespace[2pt]
\multicolumn{5}{l}{\footnotesize Standard errors in parentheses; clustered by
  unit. \sym{*}~$p<0.10$, \sym{**}~$p<0.05$, \sym{***}~$p<0.01$.} \\
\end{tabular}
\end{table}

\section{Discussion}

The stability of the interaction coefficient across \cref{tab:spanned}'s four
columns is the substantive finding, but for a fidelity test the layout is the
object of interest. A converter that flattens the two header levels would merge
the group titles into the specification row and lose the \verb|\cmidrule|
spans; one that ignores \verb|\addlinespace| would collapse each standard error
onto its coefficient with no visual separation; and one that strips the trailing
\verb|\multicolumn| note row would discard the star legend that gives the
superscripts their meaning. Each of these is a distinct failure mode, which is
why the fixture combines them in a single realistic table rather than testing
them in isolation.

\end{document}
